\section{3D Weak Lensing Power Spectra: Full-Sky Analysis and Spherical Decomposition}

This section provides a comprehensive treatment of weak lensing power spectra in three dimensions, emphasizing the
full-sky formalism using spin-weight spherical harmonics and spherical Bessel functions. This framework naturally
handles the geometry of observations from a single point (the observer) and incorporates the rich radial structure
accessible through photometric redshifts in modern surveys.

\subsection{Weak Lensing Fields and Their Geometric Structure}

\subsubsection{The lensing potential and its physical meaning}

Weak gravitational lensing arises from the cumulative bending of light rays due to the integrated mass distribution
along the line of sight. The fundamental quantity is the \textbf{lensing potential} $\varphi(\mathbf{r})$, which describes
the integrated effect of the gravitational potential relative to a homogeneous background:
\begin{equation}
\varphi(\mathbf{r}, \bar{r}) = 2 \int_0^{\bar{r}} d\bar{r}' \frac{\bar{r} - \bar{r}'}{\bar{r}} \Phi(\mathbf{r}', \bar{r}'),
\label{eq:lensing_potential}
\end{equation}
where $\mathbf{r}$ denotes the position on the sky (with spherical coordinates $\theta, \phi$), $\bar{r}$ is the comoving distance
to the observation shell, and $\Phi(\mathbf{r}', \bar{r}')$ is the gravitational potential perturbation at distance $\bar{r}'$.

The weighting factor $((\bar{r}-\bar{r}')/\bar{r})$ encodes the geometric lensing effect: structures close to the observer
have minimal impact (factor $\approx 1$), while structures at intermediate distances have maximum effect, and structures very far away
again have reduced impact.\footnote{This is the Born approximation geometry. The factor $(1-\bar{r}'/\bar{r})$ represents
the ``lever arm'' that determines how effectively a mass distribution at $\bar{r}'$ bends rays observed at $\bar{r}$.
Structures in the distant universe bend rays that travel through them, while structures near the observer bend already-deflected rays.}

\subsubsection{Convergence and shear: scalar and tensor projections}

From the lensing potential, two key observable quantities are derived through differential operations:

\textbf{Convergence} (the magnification field):
\begin{equation}
\kappa(\mathbf{r}) = \frac{1}{2}\nabla^2_\perp \varphi(\mathbf{r}),
\label{eq:convergence_def}
\end{equation}
where $\nabla^2_\perp$ is the Laplacian on the sphere (acting on angular coordinates only).

The convergence is a spin-0 (scalar) field that describes the isotropic magnification or demagnification of sources
due to lensing. A positive convergence causes magnification (brightness increase), while negative convergence causes
demagnification.\footnote{The factor of 1/2 is a convention. Physically, $\kappa$ measures the average projected density
along the line of sight, normalized to the critical density. Values $|\kappa| \ll 1$ ensure the Born approximation is valid.}

\textbf{Shear} (the tidal distortion field):
\begin{equation}
\gamma(\mathbf{r}) = \frac{\partial^2 \varphi}{\partial\theta^2} - \frac{1}{\sin^2\theta}\frac{\partial^2 \varphi}{\partial\phi^2},
\label{eq:shear_complex}
\end{equation}
or more precisely, in terms of Cartesian components on the tangent plane:
\begin{equation}
\gamma_1 + i\gamma_2 = \left(\frac{\partial^2}{\partial\theta^2} - \frac{1}{\sin^2\theta}\frac{\partial^2}{\partial\phi^2} - 
2i\frac{1}{\sin\theta}\frac{\partial^2}{\partial\theta\partial\phi}\right)\varphi.
\label{eq:shear_cartesian}
\end{equation}

The shear is a spin-$\pm 2$ (tensor) field describing the anisotropic tidal stretching of source images. The complex notation
$\gamma = \gamma_1 + i\gamma_2$ encodes both magnitude and orientation of the deformation.\footnote{The spin-$\pm 2$ nature arises
because shear is invariant under 180-degree rotations (it looks the same whether rotated or not), but under small angle rotations
it picks up phase factors like $e^{\pm 2i\alpha}$. This is naturally captured by spin-weight spherical harmonics.}

\subsection{Spin-Weight Spherical Harmonics Decomposition}

\subsubsection{Standard and spin-weighted spherical harmonics}

For scalar fields like convergence, the standard spherical harmonic expansion suffices:
\begin{equation}
\kappa(\theta, \phi) = \sum_{\ell=0}^\infty \sum_{m=-\ell}^\ell \kappa_{\ell m} Y_{\ell m}(\theta, \phi),
\label{eq:scalar_harmonic_expansion}
\end{equation}
where $Y_{\ell m}(\theta, \phi)$ are the standard (spin-0) spherical harmonics, and $\kappa_{\ell m}$ are the harmonic coefficients
extracted via:
\begin{equation}
\kappa_{\ell m} = \int d\Omega \, Y^*_{\ell m}(\theta, \phi) \kappa(\theta, \phi).
\label{eq:scalar_coefficients}
\end{equation}

For tensor fields like shear, we must use \textbf{spin-weight spherical harmonics} ${}_s Y_{\ell m}(\theta, \phi)$, which transform
appropriately under rotations:\footnote{A quantity with spin weight $s$ transforms as $e^{is\alpha}$ under a rotation by angle $\alpha$
about the radial direction. Shear (spin-$\pm 2$) transforms as $e^{\pm 2i\alpha}$, meaning a 90-degree rotation changes its sign.
This is the correct behavior for a quadrupole distortion pattern.}

\begin{equation}
\gamma(\mathbf{r}) = \sum_{\ell=2}^\infty \sum_{m=-\ell}^\ell \left[\gamma^E_{\ell m} {}_{-2}Y_{\ell m}(\theta, \phi) + 
\gamma^B_{\ell m} {}_{+2}Y_{\ell m}(\theta, \phi)\right],
\label{eq:shear_harmonic_expansion}
\end{equation}
where ${}_{-2}Y_{\ell m}$ and ${}_{+2}Y_{\ell m}$ are the spin-$-2$ and spin-$+2$ harmonics. The coefficients $\gamma^E_{\ell m}$ and
$\gamma^B_{\ell m}$ are the \textbf{E-mode} (curl-free, gradient part) and \textbf{B-mode} (divergence-free, curl part)
amplitudes.\footnote{Pure gravitational lensing produces only E-modes. B-modes signal either weak lensing from topological defects,
B-mode contributions from systematics, or artifacts of incomplete sky coverage. This E/B decomposition is one of the most powerful
diagnostics in weak lensing.}

\subsubsection{The relation between convergence, shear, and lensing potential}

In harmonic space, differential operators on the sphere become algebraic operations. The relations between the lensing potential
$\varphi$ and the convergence and shear are:
\begin{equation}
\kappa_{\ell m} = \frac{\ell(\ell+1)}{2} \varphi_{\ell m},
\label{eq:convergence_potential_relation}
\end{equation}
showing that convergence is related to the Laplacian (which becomes $\ell(\ell+1)$ in harmonic space) of the potential.

For shear, the relation is:
\begin{equation}
\gamma^E_{\ell m} = -\frac{(\ell-1)(\ell+2)}{2} \varphi_{\ell m},
\label{eq:shear_potential_relation}
\end{equation}
showing that shear depends on higher derivatives of the potential (fourth derivatives, giving a factor $\sim \ell^4$).

These relations are the fundamental bridge between the line-of-sight integrated gravitational potential and the observable
lensing fields.\footnote{The absence of a B-mode in Eq. \ref{eq:shear_harmonic_expansion} reflects the fact that pure lensing
from a scalar potential produces only curl-free (E-mode) shear distortions. The presence of B-modes would indicate either
systematic errors or contributions from non-scalar sources (like primordial tensor modes or topological defects).}

\subsection{3D Power Spectrum in Spherical Bessel Basis}

\subsubsection{The SFB expansion for weak lensing}

When sources have a range of redshifts (or comoving distances) accessible through photometric redshifts, we can further decompose
the fields in the radial direction using spherical Bessel functions:
\begin{equation}
\kappa(\mathbf{r}) = \sqrt{\frac{2}{\pi}} \sum_{\ell,m} \int_0^\infty dk \, \kappa_{\ell m}(k) \, k \, j_\ell(kr) \, Y_{\ell m}(\theta, \phi),
\label{eq:convergence_sfb_expansion}
\end{equation}
and similarly for shear using spin-weighted harmonics. The radial variable $r$ is the comoving distance, $k$ is the conjugate
radial wavenumber, and $j_\ell(kr)$ are the spherical Bessel functions of order $\ell$.

The SFB coefficients are extracted via:
\begin{equation}
\kappa_{\ell m}(k) = \sqrt{\frac{2}{\pi}} \int_0^\infty dr \, r^2
\int d\Omega \, \kappa(\mathbf{r}) \, k \, j_\ell(kr) \, Y^*_{\ell m}(\theta, \phi).
\label{eq:sfb_inverse_convergence}
\end{equation}

This decomposition has several advantages for weak lensing:
\begin{itemize}
\item It naturally separates angular ($\ell, m$) and radial ($k$) information.
\item Redshift information from photometric redshifts maps directly onto the radial structure.
\item The spherical Bessel functions respect boundary conditions (smoothness at origin, appropriate behavior at infinity).
\item Mode mixing from finite surveys is naturally captured by window functions.
\end{itemize}

\subsubsection{The 3D power spectrum definition}

The 3D power spectrum is defined as the two-point correlation of SFB coefficients:
\begin{equation}
\bigl\langle \kappa_{\ell m}(k) \kappa^*_{\ell' m'}(k') \bigr\rangle = C^{\kappa\kappa}_\ell(k, k') \, \delta_{\ell\ell'} \delta_{mm'},
\label{eq:3d_power_spectrum_kappa}
\end{equation}
where the observed spectrum $C^{\kappa\kappa}_\ell(k,k')$ is a matrix in wavenumber space at each multipole $\ell$.

Under the idealized condition of statistical isotropy and homogeneity (SIH), the power spectrum is independent of $\ell$ and $m$:
\begin{equation}
C_\ell(k,k') = P(k) \delta(k-k') \quad \text{(SIH limit)},
\label{eq:power_sfb_isotropic_limit}
\end{equation}
recovering a single function $P(k)$ that depends only on the radial wavenumber.

\subsubsection{Window functions and mode coupling}

For realistic surveys with finite volume and radial selection function $\phi(r)$, the observed power spectrum is related to
the true underlying spectrum through a convolution with window functions:
\begin{equation}
C^{\text{obs}}_\ell(k, k') = \left(\frac{2}{\pi}\right)^2 \int_0^\infty dk'' \, k''^2 \, P(k'') \, 
W_\ell(k, k'') \, W_\ell(k', k''),
\label{eq:observed_power_window_general}
\end{equation}
where the window function is defined as:
\begin{equation}
W_\ell(k, k') = \int_0^\infty dr \, r^2 \, \phi(r) \, j_\ell(kr) \, j_\ell(k'r).
\label{eq:window_function_sfb}
\end{equation}

This window function encodes how the survey selection function $\phi(r)$ causes mode mixing between different radial wavenumbers.
The product structure $W_\ell(k, k'') \, W_\ell(k', k'')$ shows that the convolution acts symmetrically on both observed modes.\footnote{The
factor $(2/\pi)^2$ is a normalization arising from the orthonormality of the SFB basis. The integral over $k''$ performs a convolution:
each true mode $k''$ scatters into the observed modes $k$ and $k'$ with amplitude given by the product of window functions.}

\subsection{Convergence and Shear Power Spectra}

\subsubsection{Convergence power spectrum in terms of lensing potential}

Starting from the relation $\kappa_{\ell m} = [\ell(\ell+1)/2] \varphi_{\ell m}$ in harmonic space, the convergence power spectrum
in the SFB basis becomes:
\begin{equation}
C_\ell^{\kappa\kappa}(k_1, k_2) = \frac{\ell^2(\ell+1)^2}{4} \, C_\ell^{\varphi\varphi}(k_1, k_2),
\label{eq:convergence_power_from_potential}
\end{equation}
where $C_\ell^{\varphi\varphi}(k_1, k_2)$ is the power spectrum of the lensing potential.

The factor $[\ell(\ell+1)/2]^2$ is a key feature: it means convergence power is heavily suppressed at low multipoles (large angles)
and enhanced at high multipoles (small angles). This is the mathematical expression of the fact that weak lensing is most visible
at small scales where the second derivative of the potential is large.\footnote{For low $\ell$ (large angles), spatial variations
are gradual and the second derivative is small. For high $\ell$, spatial variations are rapid and the second derivative is large,
making the lensing effect more pronounced.}

\subsubsection{Shear power spectrum and E/B modes}

The shear field involves fourth derivatives of the potential, leading to:
\begin{equation}
C_\ell^{\gamma\gamma}(k_1, k_2) = \frac{(\ell-1)(\ell+2)}{4} \frac{(\ell+1)(\ell+2)}{(\ell-1)\ell} \, C_\ell^{\varphi\varphi}(k_1, k_2),
\label{eq:shear_power_intermediate}
\end{equation}
which simplifies to:
\begin{equation}
C_\ell^{\gamma\gamma}(k_1, k_2) = \frac{(\ell-1)\ell(\ell+1)(\ell+2)}{4} \, C_\ell^{\varphi\varphi}(k_1, k_2).
\label{eq:shear_power_final}
\end{equation}

The factorial-like factor $[(\ell-1)\ell(\ell+1)(\ell+2)]$ represents the angular momentum enhancement for a spin-2 field.
It vanishes as $\ell \to 0$ and grows rapidly with $\ell$, reflecting that shear is an inherently small-angle phenomenon.\footnote{This
growth with $\ell$ explains why shear-based weak lensing measurements are sensitive to small scales and high-$\ell$ modes. Convergence,
depending on only the second derivative, is comparatively less sensitive to very small-scale structure.}

\textbf{E-mode power spectrum}:
\begin{equation}
C_\ell^{EE}(k_1, k_2) = C_\ell^{\gamma\gamma}(k_1, k_2),
\label{eq:e_mode_power}
\end{equation}
containing all the curl-free (gradient) component of the shear from gravitational lensing.

\textbf{B-mode power spectrum}:
\begin{equation}
C_\ell^{BB}(k_1, k_2) = 0 \quad \text{(from pure lensing)}.
\label{eq:b_mode_power_lensing}
\end{equation}

Any detected B-modes indicate systematics, incomplete sky coverage, or exotic physics beyond standard gravitational lensing.\footnote{In
practice, sky masks (due to Galactic dust or observational constraints) cause E-to-B leakage, generating spurious B-modes. Careful
treatment of the mask geometry is essential for interpreting measured B-modes as physical signals.}

\subsubsection{Cartesian shear components}

On the flat tangent plane around a given direction, the shear has two Cartesian components $(\gamma_1, \gamma_2)$ with equal amplitude
(by isotropy):
\begin{equation}
C_\ell^{\gamma_1\gamma_1}(k_1, k_2) = C_\ell^{\gamma_2\gamma_2}(k_1, k_2) = C_\ell^{\gamma\gamma}(k_1, k_2),
\label{eq:cartesian_shear_equal}
\end{equation}
and no cross-correlation:
\begin{equation}
C_\ell^{\gamma_1\gamma_2}(k_1, k_2) = 0.
\label{eq:cartesian_shear_uncorrelated}
\end{equation}

These relations hold by symmetry: the shear field has no preferred direction on the sky, so the two components must be equivalent.\footnote{In
real data, deviations from these symmetries indicate either systematics (e.g., PSF distortions) or that the field is not purely isotropic
(e.g., in the presence of foregrounds or large-scale structures breaking isotropy locally).}

\subsection{Connection to Matter Power Spectrum}

\subsubsection{Integration kernel relating lensing to dark matter}

The lensing potential is sourced by the matter density perturbations. In general relativity, the relationship involves the line-of-sight
integration of the matter power spectrum weighted by geometrical factors:
\begin{equation}
C_\ell^{\varphi\varphi}(k_1, k_2) = \frac{16}{c^4} \int_0^\infty dk \, k^2 \, I_\ell(k_1, k) \, I_\ell(k_2, k),
\label{eq:potential_power_from_matter}
\end{equation}
where the integration kernel $I_\ell(k_i, k)$ encodes the lensing geometry and source distribution:
\begin{equation}
I_\ell(k_i, k) = k_i \int_0^\infty d\bar{r} \, \bar{r} \, j_\ell(k_i \bar{r}) 
\int_0^{\bar{r}} d\bar{r}' \frac{\bar{r} - \bar{r}'}{\bar{r}} j_\ell(k \bar{r}') \, P_m(k; \bar{r}'),
\label{eq:integration_kernel}
\end{equation}
where $P_m(k; \bar{r}')$ is the matter power spectrum at comoving distance $\bar{r}'$.

This kernel has three key features:
\begin{itemize}
\item \textbf{Outer integral}: weights contributions at distance $\bar{r}$ by the Bessel function $j_\ell(k_i \bar{r})$, encoding how
structures at that distance project onto the observed angular scale.
\item \textbf{Inner integral}: sums up all contributions from structures between the observer and distance $\bar{r}$, weighted by
the ``lever arm'' factor $(\bar{r}-\bar{r}')/\bar{r}$ that quantifies the lensing efficiency.
\item \textbf{Matter power}: varies with redshift through both the growth factor $D(\bar{r}')$ and any redshift-dependent bias $b(\bar{r}')$.
\end{itemize}

\subsubsection{Including evolution and galaxy bias}

For galaxy weak lensing (lensing by galaxy distributions rather than all matter), we must account for bias and growth:
\begin{equation}
W_\ell^{\text{evol}}(k, k') = \int_0^\infty dr \, r^2 \, \phi(r) \, D(r) \, b(r) \, j_\ell(kr) \, j_\ell(k'r),
\label{eq:window_with_bias_growth}
\end{equation}
where:
\begin{itemize}
\item $\phi(r)$ is the radial selection function of observed sources.
\item $D(r)$ is the linear growth factor, describing how perturbations grow from early times to the observation epoch.
\item $b(r)$ is the galaxy bias, quantifying how galaxy fluctuations relate to matter fluctuations (typically $b > 1$).
\end{itemize}

The observed galaxy lensing power spectrum then becomes:
\begin{equation}
C_\ell^{\text{gal}}(k_1, k_2) = \left(\frac{2}{\pi}\right)^2 \int_0^\infty dk \, k^2 \, P_m(k) \, 
W_\ell^{\text{evol}}(k_1, k) \, W_\ell^{\text{evol}}(k_2, k),
\label{eq:galaxy_lensing_power}
\end{equation}

This shows that galaxy lensing is sensitive not just to the matter power spectrum, but also to the spatial variations of
bias and growth—information unavailable from convergence of lensing alone.\footnote{The separate evolution and bias factors make
galaxy lensing measurements more model-dependent but potentially more powerful for constraining cosmology, as they encode
information about the expansion history and structure growth rate independently.}

\subsection{Geometric Limits and Recovery of Known Results}

\subsubsection{The radialisation limit: deep survey}

For deep surveys where the radial selection extends over a large distance range, the condition:
\begin{equation}
r_0 k \gg \sqrt{2\ell(\ell+1)}
\label{eq:radialisation_condition_wl}
\end{equation}
is satisfied, where $r_0$ characterizes the survey depth.

In this limit, the window function becomes diagonal:
\begin{equation}
W_\ell(k, k') \approx \frac{\pi}{2k'} \delta(k - k'),
\label{eq:window_radial_limit_wl}
\end{equation}
and the observed spectrum approaches:
\begin{equation}
C^{\text{obs}}_\ell(k, k') \approx P(k) \delta(k - k'),
\label{eq:spectrum_radial_limit_wl}
\end{equation}

This is the deep survey limit where lensing information becomes purely radial (1D), independent of angular multipoles.
All information is compressed into the radial wavenumber $k$, making this limit analogous to standard 1D power spectrum analyses.\footnote{Physically,
the radialisation occurs because in a deep survey, the volume is dominated by distant structures that all contribute coherently
to a given radial mode. Angular modes with different $\ell$ tend to cancel, leaving only the radial signal.}

\subsubsection{The thin-shell limit: flat-sky approximation}

If the source redshift distribution is a thin shell at distance $r_*$:
\begin{equation}
\phi(r) = r_* \delta(r - r_*),
\label{eq:thin_shell_selection_wl}
\end{equation}
the window function becomes:
\begin{equation}
W_\ell(k, k') = r_*^3 j_\ell(kr_*) j_\ell(k'r_*),
\label{eq:window_thin_shell_wl}
\end{equation}

and the 3D spectrum projects onto a 2D spherical harmonic spectrum:
\begin{equation}
C^{\text{obs}}_\ell(k, k') = \frac{2}{\pi} k \, k' \, j_\ell(kr_*) j_\ell(k'r_*) C_\ell^{(2D)},
\label{eq:3d_to_2d_projection}
\end{equation}
where:
\begin{equation}
C_\ell^{(2D)} = \frac{2}{\pi} \int_0^\infty dk \, k^2 \, P(k) \left[\int_0^\infty dr \, r^2 \phi(r) j_\ell(kr)\right]^2.
\label{eq:2d_spectrum_definition}
\end{equation}

All radial information is lost; the spectrum depends only on $\ell$ and becomes diagonal in the $(k, k')$ space (up to the geometric
prefactors). This is the limit of traditional 2D weak lensing where sources are at a fixed redshift.\footnote{In this limit, the
full 3D information is compressed into 2D harmonic space. This is the regime of past weak lensing surveys (CFHT Legacy Survey,
COSMOS) where the redshift distribution was narrow. Modern surveys with photometric redshifts operate between the thin-shell
and radialisation limits, accessing the full 3D information.}

\subsubsection{Small-angle approximation and recovery of flat-sky results}

For $\ell \gg 1$ (corresponding to small angles on the sky $\theta \sim 1/\ell$), the spherical Bessel functions can be approximated
by their small-angle limits, and the spherical harmonic formalism reduces to the classical flat-sky weak lensing theory. The Bessel
function $j_\ell(x)$ oscillates with wavelength $\sim 2\pi/x$; for $\ell \gg 1$, these oscillations become very rapid, effectively
reproducing plane-wave behavior on small angular scales.\footnote{This is the mathematical basis for why flat-sky approximations work
so well for small-angle observations. The sphere locally looks like a plane, and the spherical harmonics become similar to plane waves,
recovering the flat-sky power spectrum formalism.}

\subsection{Equivalence of Spherical and Cartesian Decompositions}

\subsubsection{The plane-wave expansion formula}

The spherical and Cartesian Fourier decompositions are related through the fundamental plane-wave expansion in spherical harmonics:
\begin{equation}
e^{i\mathbf{k} \cdot \mathbf{r}} = 4\pi \sum_{\ell,m} i^\ell j_\ell(kr) Y_{\ell m}(\theta_k, \phi_k) Y_{\ell m}^*(\theta_r, \phi_r),
\label{eq:plane_wave_expansion}
\end{equation}
where the left side represents a plane wave in Cartesian coordinates, and the right side is its expansion in terms of radial Bessel
and angular harmonic modes.

Using orthogonality relations for spherical harmonics and Bessel functions:
\begin{equation}
\int_0^\infty k^2 \, dk \, j_\ell(kr) j_\ell(k'r) = \frac{\pi}{2k^2} \delta(k - k'),
\label{eq:bessel_orthonormality}
\end{equation}
\begin{equation}
\int d\Omega \, Y_{\ell m}(\theta, \phi) Y_{\ell' m'}^*(\theta, \phi) = \delta_{\ell\ell'} \delta_{mm'},
\label{eq:harmonic_orthonormality}
\end{equation}
one can show that a field decomposed in Cartesian Fourier space:
\begin{equation}
f(\mathbf{r}) = \frac{1}{(2\pi)^3} \int d^3\mathbf{k} \, \tilde{f}(\mathbf{k}) e^{i\mathbf{k} \cdot \mathbf{r}},
\label{eq:cartesian_decomposition_proof}
\end{equation}
yields the same two-point statistics when decomposed in the SFB basis, provided the field obeys statistical isotropy and homogeneity.

\subsubsection{Proof of equivalence for isotropic fields}

For an isotropic field, the Cartesian Fourier power spectrum is:
\begin{equation}
\langle \tilde{f}(\mathbf{k}) \tilde{f}^*(\mathbf{k}') \rangle = (2\pi)^3 P(k) \delta^3(\mathbf{k} - \mathbf{k}'),
\label{eq:cartesian_power_isotropic}
\end{equation}
where $P(k)$ depends only on the magnitude $k = |\mathbf{k}|$.

Using the plane-wave expansion and orthogonality relations, the SFB power spectrum can be shown to be:
\begin{equation}
\langle f_{\ell m}(k) f_{\ell' m'}^*(k') \rangle = C_\ell(k) \delta(k - k') \delta_{\ell\ell'} \delta_{mm'},
\label{eq:sfb_power_from_cartesian}
\end{equation}
with:
\begin{equation}
C_\ell(k) = P(k),
\label{eq:sfb_equals_cartesian}
\end{equation}
exactly matching the Cartesian Fourier result after angular integration.\footnote{This equivalence is not obvious from the definitions,
but follows from the mathematical structure of the expansions. It confirms that SFB is just another way of decomposing the same
isotropic field—the results are identical when properly integrated over all angular modes.}

\subsection{Photometric Redshift Tomography}

\subsubsection{Multi-redshift bins as 3D tomographic information}

Modern surveys provide photometric redshift estimates for each source, enabling clustering in thin redshift bins. This creates
multiple ``shells'' at different distances $r_i$, with selection functions $\phi_i(r)$ for each bin.

The cross-power spectrum between bins $i$ and $j$ is:
\begin{equation}
C_\ell^{ij}(k_1, k_2) = \left(\frac{2}{\pi}\right)^2 \int_0^\infty dk \, k^2 \, P_m(k) \,
W_\ell^i(k_1, k) \, W_\ell^j(k_2, k),
\label{eq:cross_power_tomographic}
\end{equation}
where $W_\ell^i(k_1, k)$ is the window function for bin $i$ including its selection and evolution effects.

This creates a tomographic matrix of spectra:
\begin{equation}
C_\ell^{\text{tomo}} = \begin{pmatrix}
C_\ell^{11} & C_\ell^{12} & \cdots \\
C_\ell^{21} & C_\ell^{22} & \cdots \\
\vdots & \vdots & \ddots
\end{pmatrix},
\label{eq:tomographic_matrix}
\end{equation}
capturing the full 3D lensing information with radial and angular resolution.

\subsubsection{Cosmological parameter constraints from 3D information}

The tomographic matrix contains information about:
\begin{itemize}
\item \textbf{Cosmological distances}: The positions of oscillations in $C_\ell^{ij}$ constrain $H_0$ and the expansion history.
\item \textbf{Growth of structure}: The amplitude ratios between bins at different redshifts constrain the growth factor $D(z)$.
\item \textbf{Dark energy}: Redshift-dependent evolution constrains the equation of state $w(z)$ of dark energy.
\item \textbf{Modified gravity}: Deviations from general relativity appear as non-standard coupling between the growth rate and the expansion.
\end{itemize}

The 3D treatment unlocks all of this information simultaneously, whereas 2D analyses lose the radial information and
are sensitive primarily to the late-time integrated lensing (via the Limber approximation).\footnote{The Limber approximation,
valid for high-$\ell$ modes in shallow surveys, essentially collapses the 3D information into 2D by assuming all lensing
occurs at a single effective distance. Tomographic 3D analyses go well beyond this, capturing the full redshift structure.}

\subsection{Summary and Physical Interpretation}

\begin{itemize}

\item \textbf{Weak lensing fundamentals}: The lensing potential $\varphi$, convergence $\kappa$, and shear $\gamma$ form a
hierarchical structure with increasing derivatives. Their power spectra scale as $\ell^0$, $\ell^2$, and $\ell^4$ respectively.

\item \textbf{Spin-weight decomposition}: Shear, being a spin-$\pm 2$ field, is properly decomposed into E-modes (curl-free) and
B-modes (divergence-free). Pure gravitational lensing produces only E-modes.

\item \textbf{3D SFB formalism}: The spherical Bessel decomposition naturally incorporates radial structure and redshift
information, enabling full-sky 3D analysis without flat-sky approximations.

\item \textbf{Mode coupling}: Finite surveys cause mixing between radial modes $k$ and $k'$, quantified by window functions
that depend on the survey geometry.

\item \textbf{Connection to matter}: Integration kernels relate lensing observables to the underlying matter power spectrum,
with coefficients encoding relativistic lensing geometry.

\item \textbf{Limits recover known results}: The radialisation limit (deep surveys) gives 1D radial analysis, the thin-shell
limit (narrow redshift ranges) recovers 2D harmonic analysis, and small-angle limits recover flat-sky results.

\item \textbf{Tomographic power}: Photometric redshift binning creates a matrix of cross-power spectra that captures full 3D
information and enables simultaneous constraints on distance and growth—unlocking the potential of next-generation surveys.

\end{itemize}
